% CN 6.16 Exercises (p.223--232)

% ======================================================================
\begin{exercise}{1}
    Give recurrence relations for each of the following sequences:
    \begin{enumerate}[label=(\alph*)]
        \item the positive integers;
        \item the negative even integers;
        \item the squares;
        \item the sequence whose $n$-th term is the sum of the first $n$
              positive integers;
        \item the sequence whose $n$-th term is the sum of the reciprocals of the
              first $n$ positive integers.
    \end{enumerate}
\end{exercise}

\begin{exercise}{2}
    Define $(a_n:n\in\N)$ by $a_1=2$, $a_n=2n\,a_{n-1}$.
    \begin{enumerate}[label=(\alph*)]
        \item Find the first five terms.
        \item Prove by induction on $n$ that $a_n=2^n n!$.
    \end{enumerate}
\end{exercise}

\begin{exercise}{3}
    Define $(b_n:n\in\N)$ by $b_1=3$, $b_n=b_{n-1}^2$.
    \begin{enumerate}[label=(\alph*)]
        \item Explore this sequence.
        \item Conjecture a formula for the sequence.
        \item Prove that the formula is correct, by induction.
    \end{enumerate}
\end{exercise}

\begin{exercise}{4}
    Define $(q_n:n\in\N)$ by $q_1=3$, $q_2=5$, $q_n=3q_{n-1}-2q_{n-2}$.
    \begin{enumerate}[label=(\alph*)]
        \item Explore this sequence. Find the first seven terms.
        \item Conjecture a formula for the sequence.
        \item Prove that the formula is correct, by induction.
    \end{enumerate}
\end{exercise}

% ======================================================================
\begin{exercise}{5}
    Define $(g_n:n\in\N)$ by $g_1=0$, $g_n=g_{n-1}+\log n$. Prove by induction on
    $n$ that
    \[
        g_n\le n\log n-n+H_n,
    \]
    where $H_n=1+\frac12+\cdots+\frac1n$ (see Exercise 3.11). It may help to use
    $\log(1+x)\le x$ for all $x\ge-1$, and
    \[
        \log(n-1)=\log\Bigl(n\cdot\frac{n-1}{n}\Bigr)
        =\log\Bigl(n\Bigl(1-\frac1n\Bigr)\Bigr)
        =\log n+\log\Bigl(1-\frac1n\Bigr).
    \]
    (All logarithms are natural logarithms.)
\end{exercise}

\begin{exercise}{6}
    Use the explore-conjecture-prove method to develop the best upper bound you
    can for the $n$-th term of
    \[
        r_1=1,\qquad r_n=3r_{n-1}-1,
    \]
    and prove your bound by induction. Even better, find a formula for $r_n$ and
    prove it by induction.
\end{exercise}

\begin{exercise}{7}
    The \emph{Lucas sequence} is defined by
    \[
        l_1=1,\quad l_2=3,\quad l_n=l_{n-1}+l_{n-2}\ (n\ge3).
    \]
    Adapt the approach of \S6.7.4 to develop a formula for $l_n$, and then prove
    by induction that it works.
\end{exercise}

% ======================================================================
\begin{exercise}{8}
    Define $(t_n:n\in\N)$ by $t_1=1$, $t_n=3-\frac12 t_{n-1}$.
    \begin{enumerate}[label=(\alph*)]
        \item Consider the following theorem (correct) and its ``proof''
              (incorrect). Find the errors in the proof.

              \textbf{Theorem.} For all $n\in\N$, $t_n\le\frac52$.

              \textbf{``Proof''.}
              \begin{enumerate}[label=\arabic*.]
                  \item For each $n\in\N$, let $S(n)$ be the statement that
                        $t_n\le5/2$. We must prove $S(n)$ for every $n\in\N$.
                  \item Inductive basis: Let $n=1$. Then
                        $t_1=3-\frac12\cdot1\le5/2$. So $S(1)$ holds.
                  \item Inductive step:
                  \item Assume that for some $k$ the statement $S(k)$ holds.
                        (Inductive Hypothesis.)
                  \item Consider $S(k+1)$.
                  \item $t_{k+1} = 3-\frac12 t_k$
                  \item $\phantom{t_{k+1}}\le 3-\frac12\cdot\frac52$\quad (by the
                        Inductive Hypothesis)
                  \item $\phantom{t_{k+1}}= 3-\frac54$
                  \item $\phantom{t_{k+1}}= \frac74$
                  \item $\phantom{t_{k+1}}\le\frac52$.
                  \item[11.] \hfill``$\square$''
              \end{enumerate}
        \item Write a correct proof by induction for the theorem. Two approaches:
              \begin{enumerate}[label=(\roman*)]
                  \item Prove by induction that for all $n\in\N$,
                        $1\le t_n\le\frac52$ (a stronger statement).
                  \item Prove that $S(n)$ holds for all even $n$, and then for all
                        odd $n$.
              \end{enumerate}
    \end{enumerate}
\end{exercise}

% ======================================================================
\begin{exercise}{9}
    Define $(\varphi_n:n\in\N)$ by $\varphi_1=1$,
    $\varphi_n=1+\dfrac1{\varphi_{n-1}}$.
    \begin{enumerate}[label=(\alph*)]
        \item Find the first 20 terms (spreadsheet or program).
        \item Conjecture a value for the limit $\varphi$ as $n\to\infty$.
        \item If $\varphi_n\approx\varphi_{n-1}$, the limit should satisfy
              $\varphi=1+\frac1\varphi$. Solve this to derive a conjectured exact
              value for the limit.
        \item Determine $N$ so that for all $n\ge N$,
              $|\varphi_n-\varphi|<0.001$. (Computationally is fine.)
        \item Prove that if $\varphi_n<\varphi+0.001$ then
              $\varphi_{n+1}>\varphi-0.001$.
              Hints: $w+0.001=w\bigl(1+\frac{0.001}w\bigr)$; and for all $x>0$,
              $\frac1{1+x}>1-x$.
        \item Prove that if $\varphi_n>\varphi-0.001$ then
              $\varphi_{n+1}<\varphi+0.001$.
              Hint: for $0<x<\frac13$, $\frac1{1-x}<1+\frac32x$.
        \item Deduce that if $\varphi-0.001<\varphi_n<\varphi+0.001$ then
              $\varphi-0.001<\varphi_{n+1}<\varphi+0.001$.
        \item Use (g) to prove by induction on $n$ that for all $n\ge N$,
              $\varphi-0.001<\varphi_n<\varphi+0.001$.
    \end{enumerate}
\end{exercise}

% ======================================================================
\begin{exercise}{10}
    \textbf{The Limit Game}, played on a sequence $(a_n)_{n=1}^\infty$ by Lim and
    Una:
    \begin{enumerate}[label=\arabic*.]
        \item Lim chooses a number $a$ (a proposed limit).
        \item Una chooses a real $\varepsilon>0$.
        \item Lim chooses a positive integer $N$.
        \item Una chooses a positive integer $n\ge N$.
    \end{enumerate}
    Lim wins if $|a_n-a|<\varepsilon$; Una wins if $|a_n-a|\ge\varepsilon$.

    Example: sequence $((-1)^{n+1}/n)$. Lim chooses $a=\frac12$, Una
    $\varepsilon=\frac14$, Lim $N=3$, Una $n=5$. Una wins since
    $|a_5-a|=|\frac15-\frac12|=\frac3{10}\ge\frac14$.
    \begin{enumerate}[label=(\alph*)]
        \item In this play, could Una have possibly lost on her last move?
        \item Given the first two moves as above, could Lim have chosen an $N$
              guaranteeing a win?
        \item Given Lim's first move, are all Una's options equally good? Is she
              certain to win whatever $\varepsilon$ she chooses? What advice
              would you give her?
        \item What should Lim choose on the very first move, and why?
        \item Play the game with someone, on sequences of your choice.
        \item In general, what property must a sequence have for Lim to have a
              winning strategy?
        \item Propose a sequence for which Una has a winning strategy, and
              describe it.
        \item In general, what property must a sequence have for Una to have a
              winning strategy?
    \end{enumerate}
\end{exercise}

\begin{exercise}{11}
    \textbf{The Big-O Game}, on a sequence $(a_n)$ and function
    $f:\N\to\R_0^+$, played by Oh-Yes and Oh-No:
    \begin{enumerate}[label=\arabic*.]
        \item Oh-Yes chooses a constant $c\in\R^+$ and a positive integer
              threshold $N$.
        \item Oh-No chooses a positive integer $n\ge N$.
    \end{enumerate}
    Oh-Yes wins if $a_n\le c f(n)$; Oh-No wins if $a_n>cf(n)$.

    Example: $a_n=n^2$, $f(n)=n^{1.5}$. Oh-Yes chooses $c=100$, $N=4$; Oh-No
    chooses $n=10^6$, and wins since $a_n=10^{12}>10^{11}=100\,(10^6)^{1.5}$.
    \begin{enumerate}[label=(\alph*)]
        \item In this play, could Oh-No have possibly lost on their move?
        \item What could Oh-Yes have chosen in order to win, and why?
        \item Play the game with someone.
        \item In general, what relationship must hold between the sequence and
              function for Oh-Yes to have a winning strategy?
        \item Propose a sequence and function for which Oh-No has a winning
              strategy, and describe it.
        \item In general, how must they be related for Oh-No to have a winning
              strategy?
    \end{enumerate}
\end{exercise}

% ======================================================================
\begin{exercise}{12}
    Give the simplest big-O expression you can for each of the following.
    \begin{enumerate}[label=(\alph*)]
        \item $1+4n+6n^2+4n^3+n^4$
        \item $\binom n2$
        \item $\binom n3$
        \item the sum of the first $n$ terms of an arithmetic series with first
              term 1 and common difference $d$
        \item $2^n+2^{n-1}+\cdots+2+1$
        \item $(n!)^{1/n}$
        \item $n^{1/n}$
    \end{enumerate}
\end{exercise}

\begin{exercise}{13}
    \begin{enumerate}[label=(\alph*)]
        \item How can the number $\underbrace{99\cdots99}_{n\text{ digits}}$ be
              interpreted as the sum of a finite geometric series? What are $a$,
              $r$ and $n$? What does the formula for $S_n$ give? Does this make
              sense?
        \item What about the number $0.999999\ldots$?
    \end{enumerate}
\end{exercise}

\begin{exercise}{14}
    \begin{enumerate}[label=(\alph*)]
        \item Using (6.47) $\bigl(H_n\le\log_e n+\gamma+\frac1{2n}\bigr)$, show
              that $H_n=O(\log n)$.
        \item Hence prove that
              \[
                  \frac n1+\frac n2+\frac n3+\cdots+\frac n{n-1}+\frac nn
                  = O(n\log n).
              \]
    \end{enumerate}
\end{exercise}

\begin{exercise}{15}
    Recall $(g_n)$ from Exercise 5, with closed form
    $g_n=\sum_{i=1}^n\log i$.
    \begin{enumerate}[label=(\alph*)]
        \item Use Exercise 5 and (6.47) to show that
              $g_n\le n\log n-n+\log n+c$ for some positive constant $c$.
        \item Hence show that $g_n=O(n\log n)$.
        \item Use (a) to show that $e^{g_n}\le\bigl(\frac ne\bigr)^n n\,e^c$.
        \item Prove that $e^{g_n}=n!$.
        \item Hence show that $n!=O\Bigl(\bigl(\frac ne\bigr)^n n\Bigr)$.
    \end{enumerate}
\end{exercise}
